\part{Type Theoretic Interpretations}
\label{part:XI}

Part XI interprets states, permissions, and commitments as types. Its stated payoff is the precise gap between formal inhabitation and real-world executability, and the separate roles of preservation and progress.

\begin{housing}Repair commitments are typed with deadlines and explicit failure branches, and degraded operating modes carry evidence.\end{housing}

\chapter{Economic States and Admissible Actions as Types}\label{ch:51}
Describes states, resources, permissions, and actions with specified preconditions as types. Dependent types express actions whose availability depends on current resources or established authority. Distinguishes formal inhabitation from real-world executability.
\[
\operatorname{Act}:\operatorname{State}\rightarrow\operatorname{Type}.
\]
\chspec{\Pref{P02}, \Pref{P20}.}
\scaffold

\chapter{Linear and Affine Resources}\label{ch:52}
Models resources that cannot be duplicated freely, or that may be discarded under specified rules, with linear and affine type systems. Money, physical goods, permissions, information, and obligations need different treatments, and no single resource discipline is imposed on all economic objects.
\chspec{\Pref{P20}.}
\scaffold

\chapter{Obligations, Deadlines, and Protocols}\label{ch:53}
Expresses exchanges and institutional procedures over time with session types and indexed judgments. Examines escrow, delivery, payment, default, and repair protocols, including explicit failure branches.
\chspec{\Pref{P07}, \Pref{P08}, \Pref{P20}, \Pref{P28}.}
\scaffold

\chapter{Proof-Carrying Commitments and Verified Degradation}\label{ch:54}
Commitments carry evidence of specified properties, and degraded operating modes require evidence that required functions remain supported. Small formal examples show how proofs depend on resource models, trusted observations, and assumptions about counterparties.
\chspec{\Pref{P03}, \Pref{P06}, \Pref{P14}, \Pref{P16}, \Pref{P20}.}
\draftnote{Candidate import: reconstruction by replay with canonical encoding, parent and tick validation, duplicate-operation handling, and refusal records, stated with the distinction between replay integrity and historical authority.}
\scaffold

\chapter{Refinement, Interpretation, and the Limits of Formal Assurance}\label{ch:55}
Connects abstract economic obligations with executable checks through refinement types. Distinguishes a verified implementation from a justified specification and explains why formal consistency cannot establish fairness, truthful input data, or ecological adequacy by itself.
\chspec{\Pref{P15}, \Pref{P17}.}
\scaffold
